\documentclass[11pt,a4paper]{article}
\usepackage[utf8]{inputenc}
\usepackage[margin=1in]{geometry}
\usepackage{amsmath,amssymb,amsfonts,amsthm}
\usepackage{booktabs}
\usepackage{hyperref}
\usepackage{microtype}

\hypersetup{
    colorlinks=true,
    linkcolor=blue,
    citecolor=blue,
    urlcolor=blue
}

\newtheorem{theorem}{Theorem}
\newtheorem{definition}{Definition}
\newtheorem{proposition}{Proposition}
\newtheorem{remark}{Remark}

\title{\textbf{Two Moduli Spaces and Their Topology}\\[0.5em]
\large A Synthesized Monograph Based on John Milnor's UNAM Lecture (75 A\~nos de Matem\'aticas en M\'exico)}
\author{\textbf{John Milnor}\thanks{Institute for Mathematical Sciences, Stony Brook University. Joint work with \textbf{Araceli Bonifant}. Synthesized from the inaugural keynote lecture delivered at the Auditorio Alberto Barajas Celis, Facultad de Ciencias / Instituto de Matem\'aticas, UNAM, December 4, 2017.}}
\date{December 2017}

\begin{document}

\maketitle

\begin{abstract}
This monograph presents an exposition of the inaugural keynote address \textit{``Two Moduli Spaces''} delivered by John Milnor at the Universidad Nacional Aut\'onoma de M\'exico (UNAM) on December 4, 2017, celebrating 75 years of mathematics in Mexico. Based on joint work with Araceli Bonifant, Milnor examines two fundamental algebro-geometric quotient spaces and analyzes their non-Hausdorff, non-manifold topological phenomena (``awkward topologies''). The first moduli space classifies effective zero-dimensional divisors (unordered point configurations) on the Riemann sphere $\widehat{\mathbb{C}} \cong \mathbb{P}^1$ modulo projective M\"obius transformations $\mathrm{PSL}(2, \mathbb{C})$. The second investigates the compactified space of plane algebraic curves of fixed degree $d$ in the complex projective plane $\mathbb{P}^2$ modulo projective linear automorphisms $\mathrm{PGL}(3, \mathbb{C})$. We detail the classical projective invariants, the geometric invariant theory (GIT) stability criteria (stable, semi-stable, and strictly unstable orbits), the emergence of topological non-separability, and their implications for holomorphic dynamics.
\end{abstract}

\tableofcontents
\newpage

\section{Introduction and Historical Context}

On December 4, 2017, the Instituto de Matem\'aticas and the Facultad de Ciencias of the Universidad Nacional Aut\'onoma de M\'exico (UNAM) inaugurated the international conference \textit{``75 a\~nos de matem\'aticas en M\'exico''}. The opening lecture was delivered by John Milnor (Fields Medalist 1962, Wolf Prize 1989, Abel Prize 2011). Milnor highlighted his personal connection to Mexico, having first visited Ciudad Universitaria in 1956 as a young mathematician participating in Solomon Lefschetz's renowned Algebraic Topology Symposium.

In this commemorative lecture, based on collaborative research with Araceli Bonifant, Milnor explored a recurring challenge at the interface of algebraic geometry, geometric invariant theory (GIT), and complex dynamics: how the formation of quotient spaces under non-compact continuous Lie group actions generates singular, non-Hausdorff topological spaces. Specifically, he investigated two natural geometric moduli problems:
\begin{enumerate}
    \item The moduli space of effective degree-$d$ divisors on the Riemann sphere $\mathbb{P}^1$ under the action of the M\"obius group $\mathrm{PSL}(2, \mathbb{C})$.
    \item The moduli space of degree-$d$ complex plane curves in $\mathbb{P}^2$ under the action of the projective linear group $\mathrm{PGL}(3, \mathbb{C})$.
\end{enumerate}

\section{First Moduli Space: Divisors on the Riemann Sphere}

\subsection{Effective Divisors and Symmetric Products}
Let $\widehat{\mathbb{C}} = \mathbb{P}^1(\mathbb{C})$ denote the Riemann sphere. An effective divisor of degree $d \ge 1$ on $\mathbb{P}^1$ is a formal sum:
\begin{equation}
D = \sum_{k=1}^m n_k p_k, \quad p_k \in \mathbb{P}^1, \quad n_k \in \mathbb{Z}_{>0}, \quad \sum_{k=1}^m n_k = d.
\end{equation}
Topologically and algebraically, the collection of all effective divisors of degree $d$ is canonically identified with the $d$-th symmetric product of the sphere:
\begin{equation}
\mathrm{Div}_d(\mathbb{P}^1) \cong \mathrm{Sym}^d(\mathbb{P}^1) = (\mathbb{P}^1)^d / S_d \cong \mathbb{P}^d(\mathbb{C}).
\end{equation}
This identification follows from representing any degree-$d$ divisor as the set of roots of a homogeneous binary form of degree $d$:
\begin{equation}
F(X, Y) = \sum_{j=0}^d a_j X^{d-j} Y^j, \quad [a_0 : a_1 : \dots : a_d] \in \mathbb{P}^d.
\end{equation}

\subsection{The M\"obius Action and Orbit Quotients}
The projective group $G = \mathrm{PSL}(2, \mathbb{C}) = \mathrm{Aut}(\mathbb{P}^1)$ acts on binary forms and their roots by linear changes of homogeneous coordinates:
\begin{equation}
\begin{pmatrix} X \\ Y \end{pmatrix} \mapsto \begin{pmatrix} \alpha & \beta \\ \gamma & \delta \end{pmatrix} \begin{pmatrix} X \\ Y \end{pmatrix}, \quad \alpha\delta - \beta\gamma = 1.
\end{equation}
The naive geometric moduli space is defined as the orbit space:
\begin{equation}
\mathcal{M}_d = \mathrm{Div}_d(\mathbb{P}^1) / \mathrm{PSL}(2, \mathbb{C}) \cong \mathbb{P}^d / \mathrm{PSL}(2, \mathbb{C}).
\end{equation}
Because $\mathrm{PSL}(2, \mathbb{C})$ is non-compact, the quotient topology on $\mathcal{M}_d$ exhibits non-Hausdorff pathologies whenever orbit closures intersect.

\subsection{Geometric Invariant Theory (GIT) and Stability}
Following David Mumford's Geometric Invariant Theory, points in $\mathbb{P}^d$ are categorized into stability classes based on root multiplicities:
\begin{definition}[Stability of Divisors]
Let $D \in \mathrm{Div}_d(\mathbb{P}^1)$ be a divisor with maximal point multiplicity $\mu = \max_k n_k$.
\begin{itemize}
    \item \textbf{Stable:} $\mu < d/2$. All points in the orbit have compact stabilizer subgroups; orbits are closed in the semi-stable locus.
    \item \textbf{Semi-stable:} $\mu \le d/2$.
    \item \textbf{Strictly Semi-stable:} $\mu = d/2$ (admissible only when $d$ is even).
    \item \textbf{Unstable:} $\mu > d/2$. The orbit closure contains the catastrophic degenerate orbit where all roots coalesce.
\end{itemize}
\end{definition}

\subsubsection{Low-Degree Examples}
\begin{itemize}
    \item \textbf{Degree $d = 3$:} A divisor consisting of 3 distinct points has all multiplicities equal to $1 < 3/2$. Since $\mathrm{PSL}(2, \mathbb{C})$ acts triply transitively on $\mathbb{P}^1$, any 3 distinct points can be uniquely mapped to $\{0, 1, \infty\}$. Thus, the moduli space of stable degree-3 divisors consists of a single point.
    \item \textbf{Degree $d = 4$:} For four points, the classical M\"obius cross-ratio:
    \begin{equation}
    \lambda = \frac{(z_1 - z_3)(z_2 - z_4)}{(z_1 - z_4)(z_2 - z_3)}
    \end{equation}
    yields the modular $j$-invariant:
    \begin{equation}
    j(\lambda) = 256 \frac{(\lambda^2 - \lambda + 1)^3}{\lambda^2(1 - \lambda)^2}.
    \end{equation}
    The GIT quotient $\mathcal{M}_4^{\mathrm{GIT}} \cong \mathbb{P}^1$ compactifies the affine $j$-line, where the unique strictly semi-stable orbit corresponds to pairs of double roots ($\mu = 2 = 4/2$).
\end{itemize}

\section{Second Moduli Space: Plane Curves in $\mathbb{P}^2$}

\subsection{Homogeneous Forms and the Projective Action}
The second moduli space analyzed by Milnor is the space of algebraic curves of degree $d$ in the complex projective plane $\mathbb{P}^2$. A curve $C$ is defined as the zero locus of a non-zero homogeneous polynomial $P(X, Y, Z)$ of degree $d$:
\begin{equation}
C = \{[X : Y : Z] \in \mathbb{P}^2 \mid P(X, Y, Z) = 0\}.
\end{equation}
The vector space of such polynomials has dimension $N+1 = \binom{d+2}{2}$. Thus, the parameter space of all degree-$d$ plane curves is the projective space:
\begin{equation}
\mathbb{P}^N = \mathbb{P}\left(\mathrm{Sym}^d(\mathbb{C}^3)^*\right), \quad N = \frac{d(d+3)}{2}.
\end{equation}
The group $G = \mathrm{PGL}(3, \mathbb{C})$ acts on $\mathbb{P}^N$ via linear automorphisms of $\mathbb{P}^2$. The moduli space of degree-$d$ plane curves modulo projective equivalence is:
\begin{equation}
\mathcal{C}_d = \mathbb{P}^N / \mathrm{PGL}(3, \mathbb{C}).
\end{equation}

\subsection{Classification of Singularities and Non-Hausdorff Topology}
Milnor and Bonifant demonstrate how the singularity types of plane curves dictate orbit degenerations:
\begin{itemize}
    \item \textbf{Smooth Curves ($d = 3$):} A smooth cubic curve is an elliptic curve, parametrized by its $j$-invariant. The GIT quotient yields $\mathbb{P}^1$, with the point at infinity corresponding to the nodal cubic.
    \item \textbf{Degenerate and Unstable Curves:} Curves possessing singular points of high multiplicity (such as cusps, tacnodes, or multiple line components) generate orbits whose closures contain more degenerate curves.
    \item \textbf{Topological Inseparability:} In the naive topological quotient $\mathbb{P}^N / \mathrm{PGL}(3, \mathbb{C})$, a curve with a higher singularity cannot be separated by disjoint open sets from the orbit of the less singular curve from which it degenerates, producing non-Hausdorff boundary boundaries.
\end{itemize}

\section{Connections to Holomorphic Dynamics}

Milnor establishes an explicit link between these moduli spaces and rational maps on $\widehat{\mathbb{C}}$. A rational function $f(z) = P(z)/Q(z)$ of degree $k$ has $2k - 2$ critical points, counted with multiplicity. The critical divisor $\mathrm{Crit}(f) \in \mathrm{Div}_{2k-2}(\mathbb{P}^1)$ is invariant under affine coordinate changes.

Understanding the topology of $\mathcal{M}_{2k-2}$ provides geometric coordinates for the moduli space of dynamical systems $\mathrm{rat}_k / \mathrm{PSL}(2, \mathbb{C})$. Milnor illustrates that dynamic bifurcations, parabolic degenerations, and escaping parameters in rational dynamics directly parallel orbit boundary degenerations in geometric invariant theory.

\section{Conclusion}

Through this comparative study of divisors on $\mathbb{P}^1$ and curves in $\mathbb{P}^2$, John Milnor and Araceli Bonifant elucidate how classical invariant theory, modern algebraic geometry, and complex dynamics converge. The apparent pathology of non-Hausdorff quotient topologies is transformed into a precise classification tool for geometric degenerations and dynamical limits.

\section*{References}
\begin{enumerate}
    \item Bonifant, A., \& Milnor, J. (2017). \textit{Two Moduli Spaces and Their Topology}. Stony Brook IMS Preprint Series.
    \item Milnor, J. (2017, December 4). \textit{Two Moduli Spaces} [Inaugural Keynote Lecture]. Conferencia Internacional 75 a\~nos de matem\'aticas en M\'exico, Auditorio Alberto Barajas Celis, Facultad de Ciencias / Instituto de Matem\'aticas, UNAM. Available at: \url{https://www.youtube.com/watch?v=wf8eW9qJ4-4}.
    \item Mumford, D., Fogarty, J., \& Kirwan, F. (1994). \textit{Geometric Invariant Theory} (3rd ed.). Springer-Verlag.
\end{enumerate}

\end{document}
